
% 来源：EPUB part077.xhtml
\section{生意能不能发展好关键且只取决于产品}

生意能不能发展好关键且只取决于公司的产品，无他。其他只是对快慢有影响。（2020-11-05）

% 来源：EPUB part078.xhtml
\subsection{差异化是用户需要但别人没能满足的东西}

\BookQuestion{网友：除了关注企业文化和商业模式，您会下功夫去研究公司产品吗？}

关注企业应该只能从关注产品开始的吧？没想到例外。（2020-10-27）

\BookQuestion{网友：那些商业模式很强大的公司，本质上是不是因为拥有强大的产品，所以才能形成强大的商业模式，比如茅台、苹果、腾讯？}

强大产品是必要条件。（2020-10-30）

\BookQuestion{网友：好的生意模式是否可以说是做出差异化的生意模式？}

倒过来讲可能好点：没有差异化产品的商业模式基本不是好的商业模式。所以投资要尽量避开产品很难长期做出差异化的公司，比如航空公司，比如太阳能组件公司。（2015-04-20）

\BookQuestion{网友：差异化是否是好的生意模式的必要条件？}

很难说必要条件，比如某些垄断企业，如中石油。（2013-03-31）

\BookQuestion{网友：做产品是不是以我为主，做自己认为对的事情就可以了。}

做产品哪敢不看需求？看不看得对是水平问题，看不看可是路线问题啊。（2010-04-03）

\BookQuestion{网友：Macbook Pro真是让人失望啊，苹果的文化貌似在笔记本上一直是停滞的。}

产品最重要的是有些人喜欢而不是所有人都喜欢的。（2016-10-30）

\BookQuestion{网友：平板电脑的细分市场一直想不明白。iPad的火爆是针对平板电脑市场内的细分，还是相对于原有PC市场开创一个新的品类？目前在规划产品时，还是没有办法跳出竞品导向的思维方式，对于这类新生的产品形态，该用什么样的方式进入？我对“敢为天下后”又产生了一些困惑。}

差异化指的就是用户需求满意度上的差异，绝对不仅仅指的是外观。产品的差异化不是指所谓的“与众不同”，而是指“与众不同”的东西正好是用户需要而其他人没能够满足的东西。当一个产品找到的差异化正好是很多很多用户需要的东西时，那这个产品大概就很成功了。

差异化的东西是在不断地变化的，当大家（很多公司）都有了的差异化就会变成基本需求。有时候好的产品的差异化东西不一定需要很多，有时候哪怕有一个也会让公司（或产品）很成功。

比如，如果你能想清楚麦当劳的差异化是什么，也许你就会明白些。至于如何找到差异化的问题，那是如何把事情做对的范畴的事情，那可不是一朝一夕之“功夫”啊。苹果能有现在这些个好东西，那也是十几二十年积累的“厚积薄发”而已。（2011-01-10）

\BookQuestion{网友：能否讲讲您眼中的差异化，小公司该如何做好差异化？}

你自己眼里有差异化吗？如果没有，别人讲没用，如果你懂什么是差异化其实也不用别人说，呵呵。（2012-02-21）

明白差异化最简单的办法就是把自己当消费者，想象一下自己的体验，自己常去哪里，因为什么而改变等等。（2012-05-18）

\BookQuestion{网友：步步高和OPPO进入的市场看起来鸡肋，但是步步高和OPPO都能提供不错的产品和创造价值。显然，这两家公司对进入新市场有一个相当严谨而且理性的判断。你们如何做出进入新市场的决策？}

智能手机市场算鸡肋吗？鸡肋市场的意思是指大公司食之无味弃之可惜的市场，比如学习机。实际上是个细分市场的概念，说通了就是差异化的意思。（2013-06-29）

\BookQuestion{网友：如果云南创可贴让伤口痊愈更快，粘一些也无碍。}

你的更快是和不贴比，不是和别的产品比。产品的基本功能不能满足的话，根本就不应该上市。但产品的竞争力往往是基本功能以外的东西。就像好卖的手表不是因为更准一样。我就是昨天刚刚用过以后觉得不好才去7-11买的邦迪。（2010-08-17）

\BookQuestion{网友：差异化要实际操作下来感觉难度比较大，公司投入大，工厂配合度要很高，市场还不一定认可，风险比较高。还不如直接抄袭爆款稍微改改。特别是在价格战面前，花费很大精力做的一点差异化都被秒成渣渣。目前公司就十来个人，所以我们现在都不追求差异化了。}

不是被秒成渣渣的吧？

\BookQuestion{网友：大部分“差异化”本身就是渣渣，并不是“用户需要但同行满足不了的东西”。真正的差异化大部分都是厚积薄发的结果。小公司得想办法先活下来，国内大部分成功的公司早期都是从模仿成功产品起步的，产品差异化是发展的方向，需要功夫，只能慢慢来。}

是啊，所谓“秒成”渣渣的实际上本来就是渣渣。差异化不是简单的不同，而是用户需要但同行还没满足的东西，潜在前提是至少要有同行已经满足了用户需要的东西。苹果这样的公司搞个有差异化的产品可能需要8年10年，我们凭什么就觉得自己想搞就能很快搞出来呢？（2023-06-28）

% 来源：EPUB part079.xhtml
\subsection{没有差异化，最后就是价格战}

产品要有差异化，不然最后就是惨烈的价格竞争。（2012-07-27）

\BookQuestion{网友：企业可以通过三方面获得竞争优势：低成本、差异化和专一化。}

长期来讲，低成本是靠不住的，差异化最重要，专一是产生差异化的基本点。（2011-08-03）

\BookQuestion{网友：《穷查理宝典》里提到，一些竞争激烈的行业比如麦片，经常做一些促销打折活动，但是大家都活得很好；而有一些行业最后只剩下几个寡头，但是谁都不赚钱。我没有想通，为什么会存在这两种截然不同的现象。}

我觉得决定因素是差异化。比如航空公司，由于产品几乎没办法差异化，最后只能靠价钱。你可以查查从北京到广州的各航空公司的票价，我猜一定几乎是一样的。而麦片毕竟是吃的东西，不同牌子的东西口味不一样，买的人不会因为5\%的折扣就换口味。

iPhone也是类似的情况，习惯用iPhone的人是不大可能因为别的手机便宜一点而换的。以前人们常换手机，是因为各手机之间的差异化其实非常小，所以新的功能和外观会吸引人们去换。iPhone有可能会是实际上最便宜的手机，你只要看看你抽屉里那些已经不用了的手机就明白了。

网络游戏其实也是一个有相当差异化的产品，虽然不玩的人看起来都差不多，但对在游戏里面的中坚玩家而言，社区和感觉的差异实际上是巨大的，不那么容易换。能够持续有用户喜欢的差异化产品的公司往往就有了很好的生意模式。（2012-05-02）

产品差异化程度越低，行业内的企业越难赚到钱。航空公司是极致，就是因为产品差异化小。还有一个极致的例子，就是现在很多做太阳能硅片的，下场会比航空公司还惨。（2012-07-11）

\BookQuestion{网友：造纸、纺织是不是和硅片一样？在我看来产品也没有太大差异化。饮用水或医药算不算没有差异化的行业？}

感觉最没有差异化的产品就是硅片，因为用户最后只会关心每度电的成本，所以这个行业实际上是效率的竞争。长期而言，只有某些能真的做出高效率光电转换的产品的企业才能生存下来，而且它们的产品以后很可能就是发电厂。

航空公司的差异化也非常小，除非某航空公司的安全性被怀疑，否则消费者的关心点会主要集中在从a点到b点的成本。当然，服务态度会有一点点差异（如果服务态度很不好则会有很大差异）。所以以上两个行业的价格竞争会异常激烈。

你说的其他几个行业的差异化程度都比这两个行业或多或少要高一些。有些行业看起来差异化很小，其实有些“小东西”会造成很大的差异化，比如零售行业的地理位置，往往会让人们觉得方便比价钱重要，尤其是价钱不贵的东西。

吃喝的东西心理上的差异化其实很大的，当然渠道的差异化也很厉害，比如可口可乐的护城河中，“在哪里都能买到”这点是非常厉害的。（2012-07-13）

\BookQuestion{网友：汽车、手机、电脑、电视机这些行业中的公司绝大部分最后都会失败，幸存下来的几个进化成了最有效率的拼装工，拼装工是赚不了几个钱的。手机行业的下场不会比电脑和电视机行业好多少。}

汽车有丰田，电脑有Mac。我见过真正没什么差异化的产品很少，比如航空业、光伏、纺织等。（2013-04-06）

突然想到巴菲特说过的那句话：如果你在坑里，至少别再往下挖了。产品完全没有差异化的公司，将来早晚都会出问题的。航空公司是非常好的例子，其他类似50步和100步差别的就不细说了。（2015-12-02）

\BookQuestion{网友：不同的白酒口感不一样，喜欢喝茅台的不会因为茅台涨价而买五粮液，因为有用户黏性。家电是大众消费品，而且以功能取胜，用户黏性要差一点。对客户来说，谁价格便宜就买谁的。同样是家电行业，彩电就比冰箱更喜欢打价格战。20世纪90年代彩电行业老大长虹就率先打起价格战，而空调老大格力却没有这么做。}

价格战一般在产品差异化很小时容易发生。酒是有很大差异化的东西，至少我感觉如此。喜诗糖果也是有很大差异化的产品。

电子消费品，尤其是成熟的电子消费品的趋势往往是差异化越来越小，尤其在专利保护比较弱的国家里会更明显。（2011-01-14）

\BookQuestion{网友：您怎么看产品很难提价？}

很难提价的原因很可能是产品差异化不够大，一旦提价用户就会买竞争对手的产品或服务。（2020-12-07）

\BookQuestion{网友：如果某个企业为了增长的原因降价打压对手，同时对手这一年也做得确实不好，是否可以认为企业有些急功近利了？}

我认为，一般来讲，对企业而言降价就有点像核武器，能不用就别用，最后对谁都没好处。主动用核武的属于……（2011-01-05）

\BookQuestion{网友：价格战后总会优胜劣汰，胜出的大企业会有合适的利润，是否可以这样理解？}

价格战的最后结果往往是优不胜，比如汽车业和航空业。（2011-10-07）

\BookQuestion{网友：企业啥时候能开价格战这个大招？}

犯病的时候。（2014-09-30）

\BookQuestion{网友：企业要去库存，特别是出新产品的时候基本都会对现有产品进行调价。这个调价您怎么看？}

这个正常。（2014-10-01）

% 来源：EPUB part080.xhtml
\subsection{盲目创新是危险的}

用户买的从来就不是技术，也不是硬件，用户买的是好用的产品。

(2011-05-10)

\BookQuestion{网友：您怎么看技术在构建公司竞争力方面所起的作用？}

木桶的一块板，没有是不行的。（2014-06-18）

\BookQuestion{网友：在步步高的企业文化里没有提“创新”。}

我们也没有提要吃早餐、午餐及晚餐等等。（2010-09-06）

很难有企业不靠创新可以生存下来的，但创新是指在用户导向前提下的创新，而不是为了不同而不同的创新。我们公司把这叫差异化，也就是满足用户需要而别人还没有提供的东西。当然，用户需要而大家已经提供的东西我们绝对不能少。盲目创新是危险的，而消费者导向前提下的创新是企业生存的一个重要基础。（2010-03-08）

\BookQuestion{网友：什么是核心竞争力？}

最常见的误解是很多人分不清核心竞争力和核心技术的区别。（2011-01-20）

技术的优势从来就很少是护城河，形成技术的文化才是。（2013-08-17）

\BookQuestion{网友：苹果现在的团队为荣誉而战，库克是一个斯巴达式的清教徒，比乔布斯更柔软。这一点可以抵消他在创造性上的欠缺。}

第一，我不知道库克的创新能力；第二，一家公司如果需要依靠CEO的创新能力并不是一件好事。乔布斯实际上被神化了，他其实主要是建立了创新的文化而不是自己创新。（2013-02-13）

\BookQuestion{网友：高科技企业的核心竞争力在于持续的创新，而创新本身却具有不确定性，投资者如何在这个不确定性中收获确定的收益？}

按你这个理解思路，大概是没办法的。

“高科技企业的核心竞争力在于持续的创新”这句话就是错的，因为你很容易找出一家持续有新产品却完全没有竞争力的企业。（2013-02-18）

对我而言，把经济分成“传统经济”和“新经济”，以及把企业分成“传统企业”和“高新科技企业”也是瞎搞。（2012-06-26）

\BookQuestion{网友：我觉得苹果也挺有敢为天下后的意味，指纹识别早在很多电子产品上都用过，一直都不成功，我自己的笔记本上也有指纹识别，从来没用过。但是发布会上的苹果指纹识别，非常令人期待。此前4s推出Siri的时候也是把别人用不好的语音控制功能用到出彩。}

乔纳森·伊夫在Touch ID指纹功能的视频中说：Webelieve that technology is at it's very best, at its most empowering, when it simply disappears（我们认为，当科技完全融入生活而让人几乎意识不到它的存在时，科技就达到了最佳状态，也最具赋能作用）。（2013-09-11）

我对伊夫这句话的理解是：最牛的技术就是用户看不到的技术（技术消失的时候）。想想大家以前见过的PC上那个巨丑的指纹识别器就明白这是啥意思了。（2013-09-15）

\BookQuestion{网友：你是怎么找到你们公司产品的使命的？}

产品角度，是慢慢摸索出来的，发现不对，赶紧停。比如苹果的充电器，说了一年了，今年没推出来。没推出来，肯定是有问题没解决。没解决就不推。（2018-09-30）

\BookQuestion{网友：您在步步高时确定的是消费者导向还是技术导向？}

产品的设计当然是消费者导向，技术是隐藏在产品里的，是为产品服务的。没有很强的技术是无法实现消费者导向的。（2010-04-23）

\BookQuestion{网友：在您企业发展早期，研发投入的比例有多少？这个比例会越来越高吗？}

你是说茅台吗？可以查下他们的报表。当然，上市公司很多，你可以多看几个公司的统计结果就明白或糊涂了。

其实乔布斯对研发比例有个说法很好，你可以找来看看？（2021-12-15）

\BookQuestion{网友：我们公司正在接受上市辅导，券商的一个重点工作是帮我们把研发费用占比提高到“行业合理水平”，但是那个合理水平真的很不合理。上市有时候让人看明白／糊涂了，要做一些不合理的事情，好让自己看起来合理一些。}

为了上市提高费用比例是非常滑稽的事情，目的其实就是要给买股票的人造点假象，本身也是不诚实的行为。（2021-12-06）

\BookQuestion{网友：为什么乔布斯缔造了数字帝国，史玉柱仅是一个商业神话？}

不懂问题是什么。我觉得乔布斯如果在中国的话，应该不如史玉柱。（2010-03-26）

\BookQuestion{网友：美国近几十年来确实是技术兴国，我所在行业里就有Adobe这样由两位数学家创立的图文软件巨头。}

技术兴国是现象。（2010-03-31）

\BookQuestion{网友：中国在未来十年内也出不了一个乔布斯，出不了一个苹果！}

中国能出个马云就不错了。中国的大环境估计相当长的时间里是出不了的，这个时间会比大家想象的长很多。（2011-10-10）

不知道大家对“乔布斯”的定义是什么，但我觉得中国很长时间里还不具备出“乔布斯”的大环境。不过，中国这些年确实还是有很大进步的。当然，如果我能做“乔布斯”的话，我是不怕“鸭梨”的，大不了和苹果一起吃。

\BookQuestion{网友：您过去和现在同时在做“乔布斯”和“巴菲特”，但真不知道未来您更想做谁？}

其实我一直在做自己。真实的乔布斯并不那么美妙的，我可没兴趣当。（2011-11-28）

创新也是可以学的。最重要的一点就是如何创造自由思想的氛围。不过，在我们老中的文化氛围里做到这点确实特别难。（2012-04-27）

\BookQuestion{网友：不保护知识产权会引发恶果。}

“不保护知识产权的恶果”确实是最关键点啊。看上去保护知识产权是件重要但不那么紧急的事，现在看起来是越来越紧急了。（2012-05-12）

如果我们解决不了专利问题，我们就没办法让自己的企业创新，我们就势必长期落后于人。（2011-02-04）

专利是所有中国企业的痛啊，避无可避，尤其是走向国际时。靠自己的专利积累去防御还需要很多很多年。（2011-11-27）

\BookQuestion{网友：OPPO确实为中国电子企业争光了。在美国真的能建起来研发团队吗？这种研发模式对于主力部队几乎全在中国的公司来说运作效率如何？}

我们在美国的研发团队人不多，和国内的有很强的互补作用。另外，我们运气确实好，居然在硅谷能找到很顶尖的人加入我们这种小公司，大家工作的效率很高啊。当然，我们要走的路还很长，现在才刚刚开始。

\BookQuestion{网友：OPPO在技术上面是以整合为主，还是以研发为主？目前OPPO的视听领域有已经领先的核心专利技术吗？}

呵呵，缺什么什么重要啊。我们有挺多不错的技术，研发的投入也挺大的。恕我不能在此说得太多了。（2010-03-22）

\BookQuestion{网友：电子的东西更新很快，要持续不断地研发，会不会很累啊？}

是的，很累但很有成就感。（2010-03-25）

\BookQuestion{网友：研发投入的资金，一般多少年后能形成生产力，产生利润啊？}

N年左右吧。（2011-12-06）

% 来源：EPUB part081.xhtml
\subsection{性价比是借口}

说追求“性价比”的公司大多是在为自己的低价找借口。长寿公司大概是不太强调性价比这个概念，老百姓心里有杆秤。好货不便宜啊。

（2017-03-25）

\BookQuestion{网友：小米刚成立时说只做发烧手机，1年后做了青春版，现在又出了更低端的红米。不懂！}

低价是条最容易的路，也是一条最难的路。（2013-08-04）

\BookQuestion{网友：从低端入手有什么利弊？}

很少有人这样做的，因为成功率太低。但如果成功的话，成本是最低的。（2013-07-02）

\BookQuestion{“联想3年，如何走出阴影。”（新闻）}

这么长的文章好难看。靠低价得到的份额实际就是阴影本身，很难走出来的。看看诺基亚就明白了。（2012-02-19）

\BookQuestion{网友：OPPO FIND5真给力，1万台降500元活动冒出35万台预约。}

说明便宜还是管用啊，可是以后怎么办？（2013-06-16）

\BookQuestion{网友：步步高推出了点读机，价格在1000多块，但我在书店的渠道看，步步高点读机周围充斥着大量杂牌，价格都比步步高低，甚至有的低了一半。点读机技术门槛不高，很容易杀进来竞争，面对这些快速和大量的跟进，又要考虑质量的成本，步步高点读机的价格怎么定啊？相对的高价能守得住吗？国内价格战很容易就把一个产品做死。}

呵呵，四个轮子的都是车，价格差得比这大多了。真正的用户一般是知道差别的。比如，我们的点读机里所有地方都是原声的，都是在专业录音棚里录的。这还不包括硬件的质量差别等等。买我们的大多是因为老师推荐，这里肯定是有点道理的。

买点读机的家长一般不会为了省钱而买个像点读机的点读机，大多数人会问清楚再买的。要有合理价格最重要的是产品要有差异化，有用户想要而别人又满足不了的东西。今天我们开会时一个朋友告诉我，他非常喜欢iPad，因为It's not great for everything but extremely good for a few things（它不是样样都好，但有些方面特别棒）。（2010-04-17）

\BookQuestion{网友：消费类电子行业竞争这么激烈，为什么苹果的产品却有定价权呢？}

苹果的差异化做得好，成本控制也非常好。类似的有可口可乐、喜诗糖果，国内的茅台也类似。（2011-02-24）

\BookQuestion{网友：很多伟大的公司追求的是做消费者体验最好的产品。定位做中低端市场的公司好像很难出现伟大的企业，因为中低端市场很难在成本和消费者体验上达到一个平衡。是不是针对特定消费者，让他们感觉公司的产品或服务性价比最高，这样才是最佳的消费者体验？}

最好的产品或最好的产品体验本来就是指的一定范围内的。（2013-04-10）

\BookQuestion{网友：那苹果瞄准的是哪类客户群呢？因为苹果是财富各个层次的人都喜欢，而且4000-5000元的价格有经济基础的人们都有能力买。}

个人认为其实苹果只是努力把自己的产品做好，然后卖一个合适的价格而已。至于谁会最后落在喜欢苹果产品的范围里看起来是个运气问题。“看起来”是个运气的东西经过很多年的积累后实际上就“看起来”有点必然了。苹果有今天其实是经过30多年的积累而来的，虽然爆发的时候“看起来”有点运气。（2013-04-11）

\BookQuestion{网友：您当年在一线的时候，产品做得不好您会砸样品（肝火上升）吗？}

我什么时候砸过样机？你心目中的“老板”就是这样的？

\BookQuestion{网友：海尔张瑞敏砸过。}

他当时是不得已而为之，为的是强调品质的概念，同时也说明当时他们的品质体系还有很大问题。品质部长是有一票否决权的，没道理会为了不合格的产品打电话问老板是否可以卖。

简单讲，一个公司如果还需要靠老板砸产品来提醒员工注意品质，那他的质量体系水平就确实还处在比较初级的阶段。（2010-11-29）

质量的定义其实是一致性。（2018-07-28）

\BookQuestion{网友：曾经有批发商拿你们公司的产品然后你们自己贴了地址条，后来批发商客户找上门来了。批发商的角度当然是不希望让客户知道自己的货源。你们是悄悄放的吗？怎么衡量得失的呢？}

那时产品质量法还没出来，国家还没有硬性要求厂家必须在产品包装上打上地址。当时我们只是希望万一机器出问题时用户能找到我们，和质量法要求的意思不谋而合。（2019-06-20）

% 来源：EPUB part082.xhtml
\subsection{做企业如跳水，动作越少越好}

\BookQuestion{网友：在财经媒体上看到段先生的一句话，“做企业如跳水运动员，动作越少越好”，深受启发，也把它作为做人的原则之一。}

呵呵，我总觉得那是史玉柱说的。他老说是我说的，我自己有点不记得当时是不是用过跳水这个例子。我大概说过企业一定要聚焦，或者是焦点法则等。不过用跳水来形容非常好，谁说的不重要。（2010-03-27）

\BookQuestion{网友：聚焦是放之四海而皆准的原则吗？}

就做企业而言应该差不多吧。（2022-02-12）

多元化一般是主业不够强，想靠多元化找出路，结果找了一条岔路。（2023-07-06）

\BookQuestion{网友：为什么大公司会比小公司更容易了解些？一般大公司往往代表着复杂的业务，庞大的产品线，似乎不太容易了解，而小公司往往业务简单，产品也少，不应该是更容易了解吗？}

“大公司”就一定业务复杂吗？乔布斯笑了哈！苹果非常简单，茅台也是。小公司业务简单吗？“贾布斯”也笑了。（2017-03-21）

\BookQuestion{网友：马云说，CEO的主要任务不是寻找机会而是对机会说NO。机会太多，只能抓一个。我只能抓一只兔子，抓多了，什么都会丢掉。}

这是焦点法则。（2010-07-15）

\BookQuestion{网友：《丁磊，说不大师》这篇文章大意就是讲丁磊很保守，导致了公司仅仅游戏项目成功，其他毫无亮点，你怎么看？}

那些批评丁磊的一定不如丁磊。当然，丁磊完全可以做得更好的，他现在高尔夫打不过我，电商做不过淘宝，跑步没有刘翔快。不知道有没有人去看看与丁磊同期开始做企业的那些人里有多少比例是能活下来的？（2013-06-16）

\BookQuestion{网友：你曾说：“诺基亚聚焦在提高市场份额上了，岂有不败之。理？”（2013-02-26），现在OPPO也搞机海战术，realme每年推出的手机款式应该有20多种，您怎么看？}

如果一直如此，将来也肯定会很难受的。（2022-11-06）

真是一个蛮有意思的东西。从20世纪90年代初我就开始强调要聚焦，30年了，我们依然会走着走着就又偏了。从这里能看到苹果企业文化的强大，也能看到我们的不足。亡羊补牢未为晚也！

\BookQuestion{网友：有的空调厂家也是品类多，同样适用20平方米的空调，一个品牌对应的空调有好几个型号，差异化又不大，选择起来真困难。苹果这点就特别好，可选择少，错不了。}

看上去很简单的一件事情，做到却非常难。销售一线总是会要求品种多一点再多一点，决策者需要清醒面对这种要求。我当CEO的时候几乎每天都会有人说希望品种多一点。对我们这种产品而言，品种多一点大概是个非常愚蠢的想法。诡异的是，大家看看有多少人对此乐此不疲，包括我们自己在内。（2022-11-07）

\BookQuestion{网友：最近有一家著名的PC厂商说，今后每年要推出50款智能机器。一般来讲，大家都知道苹果这样机型越少越好，为什么有些厂商还要进行机海战术呢？}

苹果现在一年都要推两款了，说明用户的需求比较分散。推50款基本是技穷的表现，东西绝对好不了。（2013-09-16）

\BookQuestion{网友：精品战略做得好的公司比机海战略做得好的公司的资本效率和回报率要好得多，不是一般得多。}

如果大凡你能做好一样东西，你为啥还要做很多做不好的东西呢？不过，确实有些会因为做了一些好东西后就以为自己啥都能做好，于是就开始做一些自己做不好的东西了。人的精力是有限的，聚焦的人或公司最后会变得很厉害，这也是我喜欢苹果的原因。（2013-10-17）

% 来源：EPUB part083.xhtml
\subsection{敢为天下后，后中争先}

“敢为天下后”的整句话是“敢为天下后，后中争先”，没有能力后中争先的地方是绝对不该去的。（2016-03-05）

\BookQuestion{网友：最喜欢段总的一句话是“敢为人后，后发制人”。但是，未来要想做强做大，扛起行业的大旗可能还不行吧。}

给你举几个“敢为人后”的例子，苹果的iPod算是吧（之前MP3早就满大街了）？iPhone算是吧（手机不用说了）？Xbox算是吧（之前有任天堂和PlayStation等等）？PlayStation算是吧（之前有任天堂和世嘉等）？国内的例子就不举了，太多了，相信你能明白。（2010-03-26）

“对苹果来说最重要的不是抢占先机，如果你往前看，第一款MP3不是我们研发的，第一款智能手机也不是，第一款平板电脑也不是，平板其实数十年前就有了，但很少人用。但毫无疑问我们发布的第一款MP3，是初代最先进的。我们的智能手机也是，我们的平板电脑也是。对于苹果来说，我们要的不是第一，而是做到最好。”（库克）

要做就做最好！敢为天下后的苹果做到了！（2022-12-27）

\BookQuestion{网友：我觉得腾讯是一家卑鄙无耻的公司，在网上很多人都骂腾讯。}

有人骂不一定是坏事，最不好的其实是没人骂。（2010-03-28）

\BookQuestion{网友：腾讯产品很多涉及抄袭。}

不“超”谁会买呢？腾讯其实也是一个很好的“敢为人后”的好例子。（2010-03-28）

腾讯的东西有很多进化的啊！不然怎么这么多人用？为什么我们不说苹果可耻？他那个iPod不就是MP3吗？微软哪样东西不是在别人之后开始做的？不能因为自己人干得好就说人家嘛。（2010-03-30）

\BookQuestion{网友：马化腾等别人上瘾了再来收费，比尔·盖茨也这么做。我非常反感这些行为，你们都自称企业家，可是你们都体现出企业家的精神了吗？这和愚弄消费者有什么区别？}

那你觉得该怎么做？好东西让人先免费试试没什么错啊，这是最有效的广告。只有很少一部分产品可以这样做，大部分产品很难做到这一点。（2010-07-15）

\BookQuestion{网友：能讲讲"为天下先"的坏处吗？}

不知道有什么坏处，但你可以试试举几个你知道的敢为天下先到现在还不错的例子。（2010-07-16）

\BookQuestion{网友:2015年中国智能手机市场OPPO 8\% + vivo 9\% ，市场份额总共17\%。做到本分不张扬确实不简单，也验证了段大哥在《波士堂》里面说的自己公司会最终进入行业的Top前列。}

其实也简单（但绝非容易），认为自己进不了Top就不进去这个行业，另外长期进不了Top的早晚也会被淘汰的。（2015-04-28）

\BookQuestion{网友：当年步步高选产品，大多数都是当下比较热门的产品，有没担心行业发展太快导致竞争激烈产品寿命短等问题？当时怎么判断自己是否有机会？}

我们一般在开始做一个产品时总是认为这个产品是长寿的，不然我们很难做好。怎么判断就很难说清楚了，不然我们也不会有那么多失败的产品。目前看来，智能手机行业对我们蛮合适的。（2020-10-20）

\BookQuestion{网友：敢为天下后的前提是否是可以做出差异化而且客户需要的产品或服务？}

我个人理解，好的公司一般来讲就是能够持续找到用户差异化需求的那些公司吧。如果找不到差异化的东西，那“敢为天下后”的产品就会成为悲剧。所以敢为天下后的前提一定是你能够提供出你的用户群需要而别人没有或不能提供的差异化的产品。（2011-01-10）

我早年提出的敢为天下后的意思是，如果发现适合自己做的产品，即便后也是可以做的。适合这点很难。比如发现“搜索”很好但现在去做大概率是会完蛋的。（2020-11-19）

\BookQuestion{网友：如何判断能不能“后中争先”？}

简单讲，就是长远看你的综合实力是数倍于竞争对手的，或者你的产品确实有用户强烈需要但目前市场上无人能满足的东西，不然就最好别碰。（2020-11-20）

我没有具体量化的概念，毛估估觉得大概有点像“集中优势兵力打歼灭战”的意思，就是在某个时段某个局部市场你的整体力量能比对手强大三倍以上，从而最后使自己占到一席之地。这需要强大的实力，同时也需要时间，需要对的战略也需要对的战术……话虽这么说，我们自己好像从来没真的这么算过，搞得像打仗一样。

只有当你觉得一个行业或产品可以有很长久的发展时，这个做法才有意义。手机是个很好的例子，苹果也是如此，但微软就是个失败的例子。敢为天下后里能后中争先大概是个前提，就像下围棋一样，一直落后手是会输棋的。所以，敢为天下后里还有一个隐含的意思，就是：没有金刚钻不揽瓷器活。（2019-10-20）

% 来源：EPUB part084.xhtml
\subsection{案例}

\BookQuestion{网友：我所在的工业地坪涂料行业算不算没有差异化的行业？会不会像巴菲特的纺织业一样？}

看上去有点像。不过像不等于是，比如3M的产品看起来都像差异化不大的产品，实际上很多都是差异化很大的产品，比如贴纸啥的。（2012-07-16）

\BookQuestion{网友：宠物用品企业的商业模式怎么样？}

了解不细，直觉上觉得商业模式一般，产品很难有差异化。就算好吃宠物也无法表达出来。猫粮狗粮似乎都蛮标配的。（2019-10-09）

\BookQuestion{网友：您怎么看目前电视行业的格局？}

我不了解电视机行业，感觉产品难有差异化，是个苦生意。（2020-10-10）

\BookQuestion{网友：安卓阵营的手机商都成本价卖手机了，似乎小米真的开了一个坏头。如果现在大家都纷纷跟进的话，那安卓阵营最后会不会和美国的航空业一个结局？}

安卓手机不赚钱吗？小米不赚钱吗？你听谁说的？手机产品的差异化还是不小的，所以应该总是有人能够赚到钱的。航空业长期确实很难赚到钱，投资者应该尽量避开。（2015-01-11）

\BookQuestion{网友：目前手机市场同质化也很严重，但苹果手机是其中的另类。}

手机里其实有很多差异化的，就像车一样。很多人其实也觉得车的差异化很小，那是因为他们没开过别的车。

有些差异化不出事是看不出来的。不地震就“看”不出的差异。不丢手机iPhone的有些好处你就“没办法”知道。但是，有条件的用心的用户实际上是可以事先了解到的。（2013-04-12）

\BookQuestion{网友：这两年的地震，因为房子质量死了那么多人，怎么没有公司把防震作为卖点？}

呵呵，想起当年我们卖防雷击电话机的事。当时我们卖防雷击，所有人也都卖防雷击，可不到雷击没人知道谁的是真的。所以每年开春后我们的电话机销售就会有个小高潮。你这个卖点要等到地震出来后才有效？（2010-07-06）

\BookQuestion{网友：巴菲特说过要回避轮子上的公司，这里应该是指汽车行业。}

欧美走过的路表明，GDP每增长1个百分点，汽车行业就增长3个百分点。1916年的时候美国有500多家汽车厂，现在只剩下3家，中国可能会走同样的路。从上面两点来看，一方面行业的蛋糕越来越大，另一方面行业集中度越来越高。对于这个行业里的优势企业来说，应该是有机会的，为什么巴菲特选择回避呢？

这类企业的产品很难有差异化，所以长期而言很难赚钱。（2011-03-22）

\BookQuestion{“巴菲特一季度买入1000万股通用汽车。”（新闻）}

买通用汽车有点费解，生意模式可不太好啊，感觉像个雪茄烟蒂。不过通用汽车的财务数据倒是不错。

\BookQuestion{网友：为什么不太好？因为汽车是差异化不大的产品？}

差异化小是个重要方面。还有就是工会也是个麻烦，成本下不来，可能10来年后又要破产一次。如果不打算拿10年的企业为何要拿10天？非常可能这个不是巴菲特的决策，那可以看出巴菲特很宽容，但也看出巴菲特之后伯克希尔的能力确实会下降。（2012-05-18）

\BookQuestion{网友：同样是家电行业，为什么黑电、白电、厨电的净利率会相差这么大？}

应该是商业模式造成的。彩电非常难有差异化，白电和厨电产品的差异化显然是要高些的。（2017-10-18）

\BookQuestion{网友：您曾说顺丰的文化不错，您怎么看一个企业的文化好坏？}

我看到的顺丰是非常聚焦在它要做的事情上的，有利润之上的追求。（2020-12-03）

我说顺丰是家好公司的意思是，这家公司企业文化蛮好的，专注在他们想做好的事情上。但这个生意能不能有很好的利润我还搞不太懂。（2020-11-20）

\BookQuestion{网友：您如何看待快递物流行业的前景？}

行业很大但差异化不大，需要投入很多，投入期也可能很长。物流公司里，亚马逊似乎非常强大，20多年的持续投入也确实形成了自己的护城河。（2020-12-03）

\BookQuestion{网友：顺丰今天发了一季报，季度亏了大概10个亿。一直觉得顺丰的时效性是个挺宽的护城河，看来对客户而言，差异化并不够大。快递行业现在价格战非常厉害。顺带想起了Zoom，和顺丰类似。}

我依然认为顺丰和Zoom都是很好的公司，但商业模式都不是最好的那种。（2021-04-10）

我之前说过Zoom，没啥要更新的。我很欣赏Zoom的创始人和CEO，也觉得Zoom的文化很好，但就是搞不清Zoom的护城河在哪里（或者说商业模式不是很强大）。在看到强大的护城河之前，我是不敢大投的。（2022-01-30）

\BookQuestion{网友：极兔您有参与投资吗？}

确实投了，算是友情支持一下，毕竟是我们公司出去的。不过，我不喜欢物流的商业模式，不建议大家投这类公司。（2022-02-01）

\BookQuestion{网友：极兔实现单季度盈利了。}

物流是个苦生意，但也是生意。能赚钱没啥了不起，但商业模式不强大，赚的都是辛苦钱。

\BookQuestion{网友：运输行业虽然苦，但是大概率可以永续存在，物流运输业的商业模式建立在可靠稳定的需求上。}

这话说的！餐饮行业肯定是永远存在的，但大部分餐厅大概率干不了几年就换人了。（2024-08-20）

餐饮行业确实难的，起点低，管理复杂，难以保持产品的一致性……（2020-10-10）

餐饮几乎都是苦生意，高标准化的可能还不错，比如麦当劳啥的。（2024-09-02）

\BookQuestion{网友：过了两年回头看，Zoom的生意模式确实没什么护城河，长期看产品很难有差异化，感觉是个越来越难的生意。和顺丰的生意模式比起来，Zoom未来似乎更难一些。}

它们都是很难的生意。商业模式不好的生意也是生意。在有个好领导的时期，有时候苦生意也能做得还不错，但对投资而言比较难。

相对于茅台而言，它们真是难啊！（2023-01-14）

\BookQuestion{网友：饮用水算不算没有差异化的产品？我对依云矿泉水能卖这么贵一直感到很神奇。从科学的角度看它不可能有什么差异，但是被赋予故事之后又俨然很有护城河。}

人们对吃喝的东西的信任度是一个重要的差异化来源，而且这种差异不是空穴来风的。口感当然是另一个差异化的重要原因。从生意角度而言，销售渠道应该也算一种差异化吧。对卖水的企业而言，能把水卖成可口可乐那样全球都随处可见其实是相当难的。（2024-03-07）

\BookQuestion{网友：如何看待券商？}

我对券商不了解，感觉产品差异化不太大，护城河不太宽，不是我喜欢的生意。如果准入门槛非常高的话，可能也算一道门槛吧。（2020-12-02）

券商的差异化一般很小而且转换成本也很小，所以一般来说很难是特别好的投资标的。（2023-03-14）

一般来说，券商的差异化应该比航空公司还要小，转换成本非常低（换交易平台是没有成本的），价格战是难以避免的。我没有了解过任何具体的券商。（2024-01-21）

\BookQuestion{网友：如何看待电池行业？}

这类产品差异化小，变化快，建立护城河有点难度。（2019-04-29）

\BookQuestion{网友：您看得懂宁德时代吗？}

我不了解你问的公司，但你可以想想商业模式，想想产品的差异化是不是长期可以维持的，这么想可能会对理解公司有帮助。（2022-02-10）

\BookQuestion{网友：怎么看待新能源汽车以及锂电池呢？}

目前为止我还看不出谁会是最后的赢家。汽车多少还是可以做出一些差异化的，所以汽车企业可能能活得稍微好一点。电池应该很难有差异化，所以最后的结局多是价格和规模的竞争。也许三两年，也许十年八年就能看到结果了。（2023-01-15）

\BookQuestion{网友：您如何看储能？}

假设这个产品成立的前提下，直觉上储能的产品很难有差异化，技术的差异很快就会被抹平的，价格战会一直伴随这个产品或行业。（2023-03-12）

\BookQuestion{网友：搜狐的最大看点是搜狗和分拆，我认为搜狐是被低估了。}

对我而言，感觉上搜狐的最大不确定性在其对视频的不断大力投入。视频是个差异化很小的产品，一堆有钱的公司拼在一起，结局恐怕不太好。（2012-05-04）

\BookQuestion{网友：感觉视频行业很难做到差异化。但不明白为什么国内那么多公司都往里进，像百度、腾讯、搜狐、暴风、PPS等等？}

这些互联网公司大概是怕失去机会，或者是觉得不能没有。也许人们会觉得一个互联网公司如果没有视频就有差异化了（怕别人因为这个不来上网）？或者大家都是受到YouTube的鼓励？本人不在其中，不是太能想明白，但绝对不建议投资专门做视频的网站，因为非常难知道谁将会胜出，而且胜出的公司也未必就有好日子过，结局很可能会像航空公司。谁能明白航空公司这么难的生意为什么还是有那么多资本跳进去不？这个世界真奇妙啊。（2013-04-16）

\BookQuestion{网友：我看好船运公司黛安娜船舶（DSX）。行业目前处在极度困难期，但公司还是不错的，危机之前它就看到了危机，并储备了大量现金，被业内称为最保守的船东。近几年不但在危机中低成本壮大，还保持了较低的负债和一定量的现金。管理层也是租赁业中一流的！}

船运像航空公司，生意模式很不好，时间长了后，什么价都不便宜，不是我喜欢的投资。（2015-11-30）

\BookQuestion{网友：这家船舶跟别家不一样，比如高峰储备现金，低迷购买船只。}

建议你用更长的时间看。感觉你在等别人把价格抬起来以后你再买，对不？不过，我对这个行业不了解，直觉觉得这个行业的生意模式不是我喜欢的那种，要借很多钱（意味着风险大），而且产品没差异化（意味着很难有长期的好利润）。（2012-06-03）

\BookQuestion{网友：一个不会被淘汰的行业，在产能过剩时期要经历一场异常惨烈的淘汰赛。只有把那些弱者淘汰出局，行业才能恢复正常。以大起大落的航运为例，一艘日成本为2万美元的大船，在船舶紧缺时日租能达到23万美元，会吸引大量投资者投资，造成船舶过剩！当日租金跌到0.2万美金时，船舶每日都处于严重亏损。弱者无力支撑，就会破产倒闭或贱卖船舶，因为船舶维护成本很高，很多还能使用的船舶也会被拆解！当旧船被拆解到一定数量后，就会逐步平衡或紧缺，进入下一轮牛市。}

很烂的生意模式。

花了很长时间，守着一个生意模式很烂的公司，机会成本很高，心情也不好吧？花掉的时间也是机会成本。我也有过两个小时赚20多倍的记录，但不知道如何才能重复（200块到5000块）。（2015-09-09）

DSX，谁要八年前拿到现在应该蛮痛苦的。重要的话再说一遍，商业模式非常重要！（2023-10-04）

\BookQuestion{“谷歌关闭10项边缘业务，覆盖社交搜索、桌面软件和网页安全等多个领域。自从今年1月谷歌联合创始人拉里·佩奇……”（新闻）}

看来拉里·佩奇像个不错的CEO。（2011-09-03）

拉里·佩奇上台后马上就砍掉了不少没用的产品，果然厉害！（2011-10-23）

\BookQuestion{网友：雅虎中国的新闻做得真的不如新浪，也不知道马云到底能不能给整合出来，当年刚并购的时候可是雄心勃勃的，过了2-3年后他的言语中有后悔的意思，不过他好像不服气，因为他感觉自己做的企业还没失败过。}

呵呵，不服气是因为年轻啊。我觉得这可能又是个投币器的故事。人要平常心其实不容易。（2010-05-13）

\BookQuestion{网友：步步高退出生活电器，是回避红海吗？}

我们经营这么多年，退出过的产品很多，简单讲就是如果用5年10年的眼光来看这个不适合我们。（2011-04-19）

前面说过N次了，我们的核心竞争力就是我们的文化。比如决定退出生活电器领域就是一个例子，从长期来看，这个领域不适合我们，也就是说对我们而言就是个错误的方向。错了就改，不管多大的代价都是最小的代价。就这么简单，但能做到的人很少。（2011-04-20）

\BookQuestion{网友：步步高当年退出小家电市场的主要原因是什么？}

退出小家电一是因为我们不擅长，二是需要集中精力做手机。（2017-04-04）

\BookQuestion{网友：自己创业开推拿足浴店有6年，天花板比较低，想寻找新的机会。这两年有代理三家公司的产品，开店面可以维持基本收入，又有一个社交电商的加盟项目也觉得前景不错。我的能力和资源有限，是选择齐头并进还是重点突破好呢？}

我不了解你和你的生意，很难给任何建议。不过，就生意种类一般而言，做得越多，赚得越少，所以聚焦很重要。（2019-08-02）

\BookQuestion{网友：美的在做空调还有电磁炉、豆浆机等小家电，不少人说它不够专一，那步步高现在除了手机、学习机，也在做小家电。想听您谈谈对多元化的理解？}

我们这些不同大类的产品都是不同的公司做的，完全独立的不同的团队，不会为其他产品分心，所以不是多元化的产物。我不赞成一般意义上的多元化，尤其不赞成为多元化而多元化的多元化——有很多公司在经历一段不错的发展后，为了分散风险，开始搞起了多元化，结果其中有些公司很快就不用再担心风险的问题了。（2010-09-20）

\BookQuestion{网友：步步高做过VCD、DVD、电话、复读机、点读机、手机，为何说不是多元化呢？}

这都是我们三家公司在不同时期由不同团队做的产品。每个团队都应该是高度聚焦的才能做好。（2010-10-15）

\BookQuestion{网友：有些公司历来不怎么分红，后来又千方百计把利润弄没了。}

A股有些公司就有这个问题，看着利润不错，可老是不知道公司会怎么用这些钱。一看他们多元化我就有点晕。（2010-05-23）

\BookQuestion{“盛大非游戏业务盈利前景不明，拖累整体利润。”（新闻）}

没有焦点的公司早晚会有麻烦。（2011-09-01）

\BookQuestion{“盛大迷途，传媒帝国梦想渐行渐远？盛大多年来一直靠网游造血，再输血到长长的新业务线中去，以期实现从网游向传媒娱乐帝国（文学、视频、电影、支付、云计算）的转型。但随着网游收入的锐减，盛大的梦想已经渐行渐远……”（新闻）}

渐行渐远？好像一直都那么远。（2011-10-07）

\BookQuestion{“马化腾否认QPhone谣言，重申腾讯不做手机。”（新闻）}

这个我能看懂。（2012-06-13）

\BookQuestion{网友：像GE这样的成功多元化公司是非常稀少的，多元化对管理的要求能力非常高，也很容易导致大公司病。我能理解苹果在手机领域产品是高度聚焦的，但是如何理解苹果从手机、到平板、到手表，甚至有可能将来到电视和汽车领域，这样的多元化和其他公司的多元化有什么不同？}

关于苹果的东西，你可以查库克都说过啥。该说的他都说过了。你把所有的苹果发布会都看几遍大概就明白苹果是怎么回事了。如果还看不明白，我相信也没有谁能让你明白。

补充一点，理解苹果确实不容易，我似乎从来就没让不明白的任何人明白过，所以也不可能在这里几句话说明白。我的建议就是把库克说过的东西都拿出来看一下，包括每一次发布会都认真看一下，哪怕不是为了投资苹果，至少也可以看到一家好公司应该是啥样子。如果看完这些还是看不明白的话，我确实也是没办法让人明白的。（2015-02-14）

\BookQuestion{网友：苹果要是做汽车算多元化么？}

不知道多元化定义是什么，看着像是在同一个老板管理下的业务之间关联度很小的东西。苹果的各业务之间高度关联，显然不是一般所说的多元化。

\BookQuestion{网友：那腾讯有没有多元化问题？}

感觉产品上腾讯不算多元化的，他们的绝大多数产品间或多或少都是有关联或者是内在关联的。投资上似乎有点散，但仔细一想似乎也大多是有关联的。至少他们不是在追求多元化吧？（2018-11-03）

\BookQuestion{网友：我感觉不好的多元化，是这个做不好就看看换另一个事情能不能好这种搞法。}

实际情况多数可能会是因为在某个地方成功了，所以就觉得自己很厉害，想做更多的事情，不小心就跑到能力圈外面去了。有时候试一试也没啥错，所以这里其实还是蛮难划一条特别清楚的界限的。很多人的所谓多元化大概都是因为生意而去的，什么好赚钱就跟着去做点啥，和自己原来的用户群没啥关系。举个例子，如果腾讯开始做房地产了，那就叫多元化，但如果他们买块地给自己员工盖宿舍或建自己的办公楼就不算。又比如有个群众因为自己爱好开了个网站，虽然和他以前干的事情不完全一样，也不知道在不在自己的能力圈内，所以依然不算是多元化，但可以叫作创业了。（2018-11-21）

\BookQuestion{网友：史总也是多元化摔了跟头？}

史玉柱做事情还是蛮专注的，当年出问题也不是因为多元化。后来做了保健品以及游戏，在生意上也算是蛮成功的。后来这两个生意的管理团队完全不同，所以不算多元化。（2018-11-04）

\BookQuestion{网友：是不是在新的领域，自己的能力能够满足现有客户的需求，并能创造价值，获得收益，这叫多元化成功？}

多元化的意思就是俗语“不要吊死在一棵树上”，要多找几棵树。

(2018-11-21)

我看到的多元化多是因为自己在某个地方成功了，想要如法炮制到别的地方。比方说篮球打好了就跑去打棒球，或者高尔夫啥的？（这句话是开我喜欢的两个篮球运动员的玩笑哈。）（2019-01-09）

喜欢打棒球那个是乔丹，居然中间去打了一年职业比赛，然后回来又拿了三届冠军。库里和乔丹一样，在篮球以外的运动上都不如其他职业球员。（2019-01-10）

\BookQuestion{网友：请问您对TCL这些年无论在经营范围或在经营地域都实行遍地开花的做法是怎么看的？说白了，他们似乎也一直在努力，怎么利润回报就不理想呢？}

“一直在努力”就一定有利润？先要做对的事情，努力才有用。（2010-10-14）

\BookQuestion{网友：对企业来说做对的事情如何判断？譬如吉利收购沃尔沃，联想收购IBM个人电脑部门等。}

简单的东西不容易啊。做对的事情，把事情做对，就是简单不容易的典型代表。其实从外人的角度来看，我无法判断吉利的收购案或联想的收购案对他们来讲是否是件对的事情，但他们最后自己都会知道的。

关于并购，我可以讲一个我的简单理解。如果当有人本着“大不一定强，不大则一定不强。所以要做强则先做大”的想法去并购的话，那结局一定是很难看的。韦尔奇的自传里写过一些关于并购的原则，你有兴趣可以去看看。（2011-01-26）

\BookQuestion{“雅虎高管称今年将有重大收购。”（新闻）}

一看到雅虎说要开始加大并购了我就有点汗。美国雅虎的现金流还是非常不错的，但市场地位在未来会不断下降。目前雅虎的价钱里基本上就没有考虑美国雅虎的价值，这多少有点不公平但也未必一点没道理。如果雅虎未来几年为了自己“东山再起”的雄心而进行大规模收购的话，那美国雅虎就真没什么价值了。不过，到目前为止雅虎看起来还是很理性的。（2011-04-08）

\BookQuestion{“苹果CFO称谷歌125亿美元收购摩托罗拉移动价格偏高。”（新闻）}

别人谁买摩托罗拉都很难有好下场，但谷歌出手则可能有点不同。一方面说明谷歌确实急了；另一方面，谷歌的文化确实强大，或许有机会？无论如何，谷歌的文化和摩托罗拉的差异实在是太大了，如果谷歌能够搞定摩托罗拉那就是谷歌的又一个奇迹。3-5年后就能看到结果。（2011-08-19）

\BookQuestion{网友：您觉得并购摩托罗拉的决策做得怎样？}

谷歌买摩托罗拉肯定不是为了买摩托罗拉的生意，所以我无法评价。如果是为了生意而去那就是脑袋坏了。（2012-01-24）

\BookQuestion{网友：如何看苹果的收购？}

苹果在建立生态中也许可以通过买这种小公司买到一些人才和技术以及专利，以及别的我不太清楚的东西。苹果肯定不会为了买点营业额或者利润啥的去收购的。（2024-01-06）

\BookQuestion{网友：如何看待分拆？}

分拆一般是没有意义的。我能想出来的意义大概是不同部门都想通过自己的业绩来体现自己的价值，而老板又没办法协调，大概就只能这么做。（2010-04-29）

\BookQuestion{网友：对集团公司下属子公司分拆上市的股票，我本能地比较反感，总觉得可能有利益输送。段大哥怎么看？}

一般的情况下，我也不喜欢。这种情况下，如果我对创始人或者管理层不信任的话，我是不会碰的。（2010-05-25）

\BookQuestion{网友：腾讯旗下这么多生意，为什么不学搜狐和盛大把旗下有潜力的业务分拆上市呢？}

把海陆空三军分拆上市会是啥子结果？（2011-10-16）

\BookQuestion{网友：苹果分拆，是否更利于旗下的子产品部发展壮大？}

不知道你是不是觉得汽车的4个轮子拆开来跑，车会跑得更快？（2013-02-21）

\BookQuestion{网友：巴菲特的收购例子中，没有提到业绩对赌之类的。}

对赌双方一般都比较短视。作为投资方，对赌总是立于不败之地的，对接受投资的人而言，参与对赌是匪夷所思的。（2017-01-07）
